\section{Zusammenfassung}
Ziel dieser Arbeit war die Entwicklung eines automatisierbaren Montagekonzepts für SLC-Linearführungen, bei dem das Montieren und Vorspannen der Linearführung unter kontrollierten Kraft-, Positions- und Drehmomentbedingungen erfolgt. Ein besonderer Schwerpunkt lag auf der reproduzierbaren Einstellung der Vorspannung, da diese einen wesentlichen Einfluss auf die Steifigkeit, die Positioniergenauigkeit sowie das Bewegungsverhalten der Positionierer besitzt. Darüber hinaus sollte das Montagekonzept variantenfähig ausgelegt werden, um unterschiedliche SLC-Baugrößen abbilden zu können und eine spätere Prozessüberwachung sowie Automatisierung zu ermöglichen.

Auf Grundlage einer Anforderungsanalyse wurde ein modular aufgebautes Montagesystem konzipiert. Der Aufbau besteht aus einem kartesischen XYZ-Achssystem zur Positionierung der Funktionseinheiten, integrierten Schraubereinheiten für den kontrollierten Verschraubprozess sowie einer mechanischen Vorspann- und Einpressmechanik. Die verwendeten Linearachsen übernehmen ausschließlich Positionieraufgaben und sind auf eine hohe Wiederholgenauigkeit ausgelegt. Der Schraubprozess ist so ausgeführt, dass definierte Drehmomente reproduzierbar aufgebracht werden können. Durch den modularen und klar strukturierten Systemaufbau ist das Montagekonzept grundsätzlich für eine schrittweise Erhöhung des Automatisierungsgrads sowie für eine spätere Integration in industrielle Fertigungsumgebungen geeignet.

Zur Überprüfung der Funktionalität und der Einhaltung der Anforderungen wurden umfangreiche Validierungstests durchgeführt. Die Validierung erfolgte manuell in Form von Einzeltests aller relevanten Aktoren, Sensoren und Schraubereinheiten. Hierzu wurden Referenzfahrten der Achssysteme durchgeführt, Endlagen angefahren sowie wiederholte Positionierbewegungen zwischen definierten Positionen ausgeführt. Die Positioniergenauigkeit wurde mithilfe geeigneter Messmittel überprüft. Die Ergebnisse zeigen eine hinreichend höhe Wiederholgenauigkeit der Achssysteme mit einer Abweichung von 4 bis \qty{8}{\micro\meter} pro Achse. Sämtliche eingesetzten Sensoren lieferten zuverlässige Signalrückmeldungen, und alle Aktoren sowie Schraubereinheiten funktionierten innerhalb der vorgesehenen Betriebsgrenzen. Die durchgeführten Tests bestätigen die grundsätzliche Funktionsfähigkeit und Praxistauglichkeit des entwickelten Montagekonzepts.

Aufgrund des zeitlich begrenzten Rahmens der Bachelorarbeit wurde auf die vollständige Umsetzung einer softwareseitigen Automatisierung verzichtet. Insbesondere die Programmierung der Steuerungslogik, die strukturierte Kabelführung, der Aufbau eines vollständigen Schaltschrankkonzepts sowie die sicherheitstechnische Auslegung des Gesamtsystems konnten nicht abschließend realisiert werden. Diese Punkte wurden bewusst als Abgrenzung der Arbeit definiert und stellen Ansatzpunkte für weiterführende Arbeiten dar.

Als Ausblick bietet sich die vollständige Automatisierung des Montageprozesses durch die Implementierung der Steuerungssoftware an. Ergänzend sollen sicherheitstechnische Aspekte gemäß geltender Normen umgesetzt, eine kontinuierliche Prozessdatenerfassung zur Qualitätsüberwachung integriert sowie das System in eine bestehende Produktionsumgebung eingebunden werden. Das entwickelte Montagekonzept stellt somit eine belastbare Grundlage für die weitere Automatisierung des Aufbaus von SLC-Linearführungen dar.
